\documentclass[11pt,a4paper]{article}
\usepackage[margin=25mm]{geometry}
\usepackage{amsmath,amssymb,amsthm,hyperref}
\setlength{\emergencystretch}{2em}
\newtheorem{lemma}{Lemma}
\newtheorem{theorem}[lemma]{Theorem}
\newcommand{\gen}[1]{\langle#1\rangle}
\title{Saturated orbit forests and the even-degree boundary}
\author{Human Compression project: a component of the L019 recovery}
\date{}
\begin{document}\maketitle
Throughout, $G=\gen{\rho_0,\ldots,\rho_{r-1}}$ is transitive
on $n$ points, its involutions satisfy the string relations,
its noncommutation diagram is connected, and all maximal
omissions are intransitive. A colour is eligible when its edge
joins different orbits of its actual maximal omission.
The acquired unique-crossing lemma identifies a colour with
one eligible edge as a bridge between precisely two omission
orbits. A perfect split has all earlier marks supported on one
side and all later marks supported on the other.

\begin{lemma}[A saturated transitivity forest]\label{maximal}
Let $s_1,\ldots,s_t$ be an inclusion-minimal family of
involutions transitive on each of two disjoint sets of sizes
$m_1,m_2$. If $t=m_1+m_2-2$, each $s_l$ is a single
transposition on one of the sets, and the corresponding edges
form a tree on each set.
\end{lemma}
\begin{proof}
For each member choose an edge between two orbits obtained
by deleting that member. As in the acquired minimal-transitivity
proof, the selected edges form a forest. The assumed edge count
forces exactly two components, the two sets in the statement.
For any further transposition of $s_t$, join its endpoints
by the unique selected path. If this path contains an edge
of some different colour $l$, the further $t$-edge and the
rest of the path give an alternative path between that
selected $l$-edge's endpoints, using no colour $l$.
This contradicts its eligibility under deletion of $s_l$.
The path can therefore contain only the selected $t$-edge;
the further edge is that same edge. This applies to every
member and every one of its transpositions.
\end{proof}

\begin{theorem}[Saturation on two imprimitive sides]\label{sat}
Suppose a unique $i$-edge joins $a\in A$ to $b\in B$,
and $G_i$ preserves respectively $m_1,m_2\ge2$ blocks of
sizes $k_1,k_2\ge2$ on these two orbits. If $2r=n+2$,
then this cut is perfect.
\end{theorem}
\begin{proof}
Choose an inclusion-minimal subfamily $M$ of the omission
marks transitive on both block sets. The full orbit-budget
proof gives, with $c_j$ its number of point orbits on side $j$,
\[
 |M|\le m_1+m_2-2,\qquad c_j\le k_j,\qquad
 r-1-|M|\le c_1+c_2-2.
\]
Together with
$m_1+m_2+k_1+k_2-3\le(n+2)/2=r$, these force equality
in every displayed inequality. Thus $|M|=m_1+m_2-2$,
$c_j=k_j$, and there are $k_1+k_2-2$ remaining omission marks.

Lemma~\ref{maximal} applies to the block images of $M$.
The two block trees are paths: three incident transpositions
would pairwise fail to commute, whereas three distinct
integer labels cannot be pairwise consecutive. Consecutive
edges of each path have consecutive labels; the direction
cannot change without repeating a label. The block containing
$a$ is an endpoint of its path. Otherwise its two incident
labels would both belong to $\{i-1,i+1\}$, because all more
distant marks fix $a$. These two labels are distant and must
commute, a contradiction. Its unique incident label is
$i-1$ or $i+1$. The same reasoning applies at $b$, and these
two incident labels are different because each member of $M$
acts on just one block set. After reversing the order if
necessary, the two block paths therefore have labels
\[
 i-m_1+1,\ldots,i-1\quad\hbox{on }A,
 \qquad i+1,\ldots,i+m_2-1\quad\hbox{on }B,
\]
with the crossing block at the endpoint next to $i$.
These are all the members of $M$.

Each point orbit of $\gen M$ meets every block of its side.
Since there are exactly $k_j$ such orbits, each has exactly
$m_j$ points and meets each block once. Projection from any
such orbit to the block set is an equivariant bijection.
In particular the block action of $\gen M$ is faithful;
its marks act as the natural path transpositions on each
of these copies. A mark in the first block path fixes all
points of $B$, and conversely for the second path.

Let $t$ be a remaining omission mark. Its label is outside
the entire interval occupied by $M\cup\{i\}$, so it fixes
$a,b$. On either block set it centralizes the full natural
symmetric action of $M$, except possibly for the transposition
at the far endpoint of that path, if its label is adjacent
to $t$. We claim its block image is the identity in every case.
For the full path, a permutation centralizing a transitive
action and fixing one point is trivial. For the path with its
far-end edge deleted, the centralized group is the natural
$S_{m_j-1}$ on the initial $m_j-1$ blocks. When $m_j\ge3$,
its unique nontrivial orbit is preserved by every centralizer.
Fixing the crossing block and centralizing that transitive
orbit forces fixation of the whole orbit; the one remaining
block is fixed too. When $m_j=2$, fixing one of the two
blocks already forces the same conclusion. This proves the
claim, without asserting any action on the point-orbit set.

Now select one eligible actual edge for every remaining
omission mark and contract the point orbits of $\gen M$.
The acquired orbit-forest argument gives exactly two trees,
one on the $k_1$ orbits in $A$ and one on the $k_2$ orbits
in $B$. Because $t$ fixes every original block and each
$\gen M$-orbit meets that block just once, every nontrivial
transposition of $t$ joins two distinct contracted vertices.
Let $u,v$ be such a pair. If the selected path from $u$ to
$v$ contained the selected edge of a colour $l\ne t$, that
path, the actual $t$-edge, and paths within $\gen M$-orbits
would connect the endpoints of the selected $l$-edge in
$G_l$. This contradicts actual ambient eligibility. Thus
$u,v$ must be the endpoints of the one selected $t$-edge.
Every transposition of $t$ consequently lies on the same
side, $A$ or $B$.

We have proved that every omission mark acts on precisely
one side. This side cannot change at consecutive positions
on either side of $i$, since the corresponding marks would
then have disjoint supports and commute, disconnecting the
noncommutation diagram. The neighbours $i-1,i+1$ are already
in their respective block paths. Every earlier mark therefore
acts only on $A$, and every later mark only on $B$.
The split is perfect.
\end{proof}

\begin{lemma}[The two support interfaces at the boundary]\label{interfaces}
Assume $2r>n+1$ and a unique cut is not perfect. Its two sides
are nontrivial, one primitive side has its active labels in
an interval next to the cut, and after reversal $G_{<i}$
fixes that side. The acquired head-support theorem, in this
orientation, remains valid: for the last shared index $h>i$,
\[
 h<r-1,\qquad X=\operatorname{supp}(G_{>h})
 \subseteq B\setminus\{b\},\qquad 2h\le n-|X|-1,
\]
with the endpoint interchange allowed there.
\end{lemma}
\begin{proof}
The acquired orbit budget gives $2r\le n+2$ when both
sides are imprimitive. Under the present strict inequality,
its only possible equality boundary is $2r=n+2$; that cut
is perfect by Theorem~\ref{sat}. Thus a nonperfect cut
cannot have two imprimitive sides. An interior singleton
would give $2r\le n+1$, while an endpoint singleton is
already perfect. The remaining primitive-side argument is
exactly the acquired argument: disconnected active indices
force two commuting transitive factors, hence the exceptional
degree-two action of the two neighbouring involutions.
Its acquired bound is $2r\le n$. It is excluded here too.
This gives the stated interval and orientation.

The first five lemmas of the head-support manuscript have
no rank-range assumption. The fifth gives $n\ge2r$ when
a final shared mark, or a final mark on the wrong side of
a positive cut, occurs. Both are excluded by $2r>n+1$.
Connectedness then puts the final tail on a single side.
The rest of that theorem's proof is a count independent of
the rank range: two fixed path pieces at $i=0$; at $i>0$
either $|A|+1\ge2i+3$ for a transitive left side, or
$|A|+(h-i)\ge2h+1$ from its commuting-orbit and internal-pair
lemmas. All these inequalities were proved with their weak
equality cases retained. They give the claimed head bound
and its support placement under the present hypothesis.
\end{proof}

\begin{theorem}\label{boundary}
Under the standing hypotheses, $2r>n+1$ implies the existence
of a perfect split.
\end{theorem}
\begin{proof}
Use the complete acquired two-support proof with
Lemma~\ref{interfaces} supplying its orientation and head
bounds. We specify every place affected by the changed
rank assumption. The all-paired case gives $2r\le n$.
The first tail having all colours paired gives $2r\le n+1$.
The second-cut orientation uses only the last-mark bound
$n\ge2r$, and the separated-region count excludes $g\le h$
by $n\ge2r+h-g$. All are contradictions to the present
hypothesis. Thus that proof reaches exactly the same actual
sets $X,Y$ and indices $h<g$, with
\[
 X\cup Y=\Omega,\quad 2h\le n-|X|-1,\quad
 2(r-1-g)\le n-|Y|-1.
\]
If $g\ge h+2$, its middle interval is nonempty and
intransitive. Retain the full ambient interval inequality,
including its positive orbit deficit $c\ge1$:
\[
 2(g-h-1)\le |X\cap Y|-c.
\]
Summing the three bounds now gives $2r\le n+2-c\le n+1$.

If $g=h+1$, there is no middle interval. The mark $\rho_h$
is in the group defining $Y$, and $\rho_{h+1}$ is in the
group defining $X$. They fail to commute by connectedness
of the string noncommutation diagram. Their supports must
intersect; in particular $|X\cap Y|\ge1$. Sum just the
two outside bounds, retaining this overlap:
\[
 2r=2h+2(r-1-g)+4
 \le 2n-|X|-|Y|+2
 =n-|X\cap Y|+2\le n+1.
\]
This is the same contradiction and proves the theorem.
\end{proof}

\paragraph{Attribution and scope.}
The block systems and the two-support route originate in
Cameron--Fernandes--Leemans and their earlier alternating-group
work. The full actual orbit forests, paths and ambient interval
proofs are acquired premises. The saturated transitivity
argument above replaces a classification of equality diagrams:
it controls full actions on the original blocks before using
the contracted orbit forest, and never assumes that a quotient
marking has IP. The positive middle deficit and the adjacent
support overlap pay the one-unit boundary separately.

For an actual wing of a parent marking whose maximal omissions
are intransitive, every maximal omission inside that wing is
intransitive too. Indeed a transitive wing omission, together
with the other full wing and the crossing mark, would make the
corresponding parent omission transitive. Thus the theorem is
available for an even-degree symmetric wing of rank $m/2+1$,
provided its other stated hypotheses hold. It does not by itself
supply an internal split with a trivial internal factor, the
leaf-kernel conditions, or the stability bijection. Those
obligations and the complete original remain uncredited.
\end{document}
